\documentclass{article}

%% Page Margins %%
\usepackage{geometry}
\geometry{
    top = 0.75in,
    bottom = 0.75in,
    right = 0.75in,
    left = 0.75in,
}

\usepackage{amsmath}
\usepackage{graphicx}
\usepackage{parskip}

\title{Assembly Project: Dr Engels}


\author{Tyler Christopher Lin & Brandon Zento Li}

\begin{document}
\maketitle

\section{Instruction and Summary}

\begin{enumerate}

    \item Which milestones were implemented? 
    All five milestones were implemented.\\
    All features listed in this section will be numbered in accordance to the handout.\\\\
    \textbf{Easy features:}\\
    1) Implement row advancement with a timer, so that new rows of bubbles periodically push down from the top automatically.\\
    6) If the player presses the keyboard key p, display a “Paused” message on screen until they press p a second time, at which point the original game will resume.\\\\
    \textbf{Hard features:}\\
    2) Flood-fill check (horizontally and vertically) to remove all connected same-coloured bubbles after the initial match-3.\\
    3) Apply gravity so unsupported bubbles float upwards.\\
    5) If gravity causes new matches, repeat the matching and removal process until no matches remain (chain reactions).\\
    
    % TODO: List the milestone(s) and in the case of 
    %       Milestones 4 & 5, list what features you 
    %       implemented, sorted into easy and hard 
    %       categories.

    \item How to view the game:
    % TODO: specify the pixes/unit, width and height of 
    %       your game, etc.  NOTE: list these details in
    %       the header of your breakout.asm file too!
    
    \begin{enumerate}

    \item Width and height are 256 pixels, 8 units
    \item The playing board is 14 x 7, which means it is 16 x 9 with the boundary.
    \item


    \end{enumerate}

    

\begin{figure}[ht!]
    \centering
    % \includegraphics[width=0.3\textwidth]{name.png}
    \caption{caption}
    \label{Instructions}
\end{figure}

\item Game Summary:
% TODO: Tell us a little about your game.
\begin{itemize}
\item
\item
\end{itemize}

    
\end{enumerate}

\section{Attribution Table}
% TODO: If you worked in partners, tell us who was 
%       responsible for which features. Some reweighting 
%       might be possible in cases where one group member
%       deserves extra credit for the work they put in.

\begin{center}
\begin{tabular}{|| c | c ||}
\hline
 Student 1 (Name and student number) &  Student 2 (Name and student number) \\ 
 \hline
 Task & Task\\
 \hline
 Task & Task\\
 \hline
 Task & Task\\ 
 \hline
 Task & Task\\ 
 \hline
 Task & Task\\
 \hline
 Task & Task\\  
 \hline
\end{tabular}
\end{center}

% TODO: Fill out the remainder of the document as you see 
%       fit, including as much detail as you think 
%       necessary to better understand your code. 
%       You can add extra sections and subsections to 
%       help us understand why you deserve marks for 
%       features that were more challenging than they
%       might initially seem.


\end{document}\textbf{}
